\scanpage{325}\mspage{49}
% archive sheet number: 324

ens de flèches $\unsure{\mathcal{F}|F}$ qui leur donnent naissance respectivement: on notera qu'il y en a deux pour les deux triplets nommés en premier, $A,B,C$ pouvant se définir via $\mathcal{F}_{W}$ ou $\mathrm{Fibf}_{W}$ indifféremment, et de \={m} $A_{W_0},B_{W_0},C_{W_0}$ via $\mathcal{F}_{W_0}$ ou $\mathrm{Fibf}_{W_0}$, comme \inserted{déjà annoncé (4.4.3) et comme} nous le verrons plus bas \inserted{(cor. 4.4.10 pages 54-55)}. Par contre, les \unsure{triplets} \deleted{\unsure{sont}} réduits $D$ et $D^{\deleted{\unsure{'}}}_{W_0}$ sont définis chacun par \deleted{$\mathrm{Parf}_{W}$} une $\mathcal{F}$ \unsure{seult|seule}, savoir par $\mathrm{Parf}_{W}$ et $\mathrm{Parf}_{W_0}$ respectivement. (En tout, cela fait intervenir six ensembles $\mathcal{F}$ distincts, à savoir les six ensembles du diagramme (4.2.1) autres que $\mathrm{Fib}_{W}$ et $\mathrm{Fib}_{W_0}$).

L'implication non tautologique $C\overset{\bullet}{\Longrightarrow}D$ est un cas particulier de l'implication $C_{W_0}\overset{\bullet}{\Longrightarrow}D_{W_0}$ \deleted{sera \unsure{prouvée} plus bas. L'implication également non tautologique $A_{W_0}\Longrightarrow D_{W_0}$ est} \inserted{laquelle} contenue dans ce qui a été annoncé \deleted{plus \unsure{bas}} \inserted{\unsure{à l'instant|d'exprimer}}, \deleted{\unsure{i.e.}} \inserted{savoir} $\deleted{\unsure{A}}\,\unsure{C|A}_{W_0}(\mathcal{F}_{W_0})\Longrightarrow \deleted{\unsure{A}}\,\unsure{C|A}_{W_0}(\mathrm{Fibf}_{W_0})$, \deleted{\unsure{i.e.}}
\deleted{\[
 \mathcal{F}_{W_0}\subset\mathrm{Fib}_{W}\ \deleted{\unsure{\Leftarrow}}\Longrightarrow\
 \mathrm{Fibf}_{W_0}\subset\mathrm{Fib}_{W}
\]}
compte tenu que $\mathrm{Parf}_{W_0}\subset\mathrm{Fibf}_{W_0}$ \inserted{cf. (\deleted{\unsure{4}}\unsure{9|*)}\unsure{.|,}\deleted{\unsure{1)|*)}} p. 42)}. \unsure{Faisant|Evidemt}
